\chapter{Monetary Policy}

These notes follow Robert J. Shiller's Yale lecture on monetary policy, in a companion-note form curated by LazyingArt LLC. We begin with the modern picture of central banking, but the lecture very quickly refuses to leave the subject there. A central bank is not treated as a static government bureau. It is treated as a financial invention: something discovered, made to work, and then copied. From that starting point the argument moves through the Bank of England, the American imitation of that model, the widening of central banking into macroeconomic stabilization, the old reserve-requirement logic, and finally the balance-sheet mathematics of modern capital regulation.

\section{Central Banks as Financial Inventions}

The lecture opens in the present tense. Today a central bank looks like a very special government bank, the institution responsible for currency and money. Every country has paper money bearing the name of a central bank. But that definition is only the surface. We are asked almost immediately to step back and ask where such an institution came from and why nearly every country ended up with one.

This is where the lecture brings back a larger theme of the course. Financial institutions are inventions. They spread in the same way good engineering designs spread. Once somebody finds a structure that solves a real problem, the rest of the world begins to copy it. Cars converge toward certain forms, airplanes converge toward certain forms, and central banks, in the lecture's view, do the same.

That opening matters because it gives the whole chapter its motion. We are not reading a static description of a bureaucracy. We are following the history of an invention that proved useful enough to be copied widely, and useful enough that, in the recent crisis, the actions of central banks may well have prevented another Great Depression.

\section{From Goldsmith Bankers to the Bank of England}

To recover the origin, the lecture goes back to the goldsmith bankers. A depositor left gold with a goldsmith and received a paper claim. That claim could later be brought back and redeemed for specie. Once such claims began to circulate from hand to hand, paper money had effectively appeared. In its earliest modern form, then, paper money did not begin as a government invention at all. It began as a private banking practice.

That system could endure for a long time, but it had a visible weakness from the beginning. The claim was only as good as the issuer's ability and willingness to redeem it. If too many claims came back at once, the goldsmith banker might not make good on the promise.

That is the problem that leads the lecture to the Bank of England, founded in 1694. It begins simply as a bank with a special charter, but the charter matters. It gave the Bank of England a scale and a power that the smaller partnership banks did not have. It was, in the lecture's phrase, the gorilla bank. It could present the liabilities of smaller banks for payment and push them into failure.

The central point, however, is not that the large bank possessed this power. The central point is what it did with it. Instead of destroying smaller banks whenever it could, it adopted what the lecture describes as a ``live and let live'' policy under one condition: other banks had to keep deposits with it. Those deposits gave the Bank of England both leverage and a mechanism of support. If a smaller bank got into trouble, it could draw on the balance it held there; in some cases the Bank of England could lend more as well.

So the first mature central-banking logic is already in place. A dominant bank stands above the system, requires reserves to be kept with it, and turns that requirement into a stabilizing device. That is why the Bank of England matters here. It is not just the first example in a chronology. It is the institutional model that later central banks, in one form or another, copy.

\subsection{Question \& Answer}

\paragraph{Question.} Why does forcing banks to keep deposits with a dominant bank stabilize the banking system?

\paragraph{Answer.} Because the dominant bank can do two opposite things with the same power. It can destroy a weaker bank by demanding payment all at once, or it can support that bank if reserves are already lodged with it. Once the smaller banks keep deposits at the large bank, those balances become both a disciplinary device and a rescue channel. The power that could trigger collapse is converted into a lender-of-last-resort mechanism.

\section{The American Path: Suffolk, National Banking, and the Federal Reserve}

The lecture then turns to the United States, and the rhythm is important. The American story is not presented as a cleanly independent creation. It is presented as a long series of attempts to reproduce some of the advantages of the Bank of England without quite admitting that this is what we are doing.

The Suffolk Bank in Boston is the first American example. Founded in 1819, it was a private bank, but it adopted the same basic structure: New England banks were required to keep deposits with it, and the system stabilized New England against bank runs. The point is not only that Suffolk worked. The point is that it was itself an imitation of a successful financial invention.

The American currency problem in the first half of the nineteenth century then shows why such inventions mattered. Bank notes traded at discounts that depended on the issuing bank. A merchant might consult a \emph{Bank Note Reporter} and offer eighty cents on the dollar for one bank's notes and a different price for another's. The further one got from a reliable banking center, the more disorderly the system became.

The National Banking Acts of 1863 and 1864 solved one part of that problem. National banks were created, and their notes were backed through deposits with the Treasury. The result was a more uniform paper currency. National bank notes, though issued in the names of different banks, traded at par. In that sense the note problem was fixed.

But the lecture carefully separates the note problem from the banking problem. Uniform currency is not the same thing as banking stability. The United States no longer had the old discount problem in paper money, but it still had bank runs, credit booms, contractions, and banking crises. The crises of 1893 and 1907 are named in the lecture as decisive evidence that something deeper remained unfixed.

The Federal Reserve System, created in 1913 and opened in 1914, is the American answer. It is dressed in an American constitutional costume: twelve regional Reserve Banks instead of one obvious central bank, and a Board of Governors in Washington rather than a single London institution. But the underlying mechanism is familiar. Banks hold reserves as vault cash or as deposits with their Reserve Bank; if they get into trouble, the Federal Reserve can lend to them. The same stabilizing idea reappears in American form.

That is why the Fed becomes the lender of last resort. The discount window gives the mechanism a name. A bank in trouble does not simply appear and ask for money; it comes with collateral, and the central bank lends against that collateral after discounting it. Even the word ``window'' preserves the lecture's institutional imagination: one can almost picture the old teller's window and the troubled bank appearing there with its securities.

The lecture also preserves the early optimism attached to this invention. When the Federal Reserve was founded, many thought that banking crises had finally been put behind us. The British solution had at last been copied. It seemed reasonable to hope that the system would now run smoothly. The next turn of the lecture comes from showing why that hope was too early.

\section{Central Banking as Macroeconomic Regulation}

The immediate interruption is 1933. Banks had begun failing after 1929; the Federal Reserve did not stop the collapse; and just before Roosevelt took office, the American banking system fell into such disarray that the government declared a banking holiday and shut the whole system down. The lecture insists on the severity of that moment. This is not merely a date in monetary history. It is the point at which the confidence surrounding the Federal Reserve is broken.

Deposit insurance enters here as the next repair. The lecture is loose on the exact dating details, but the causal point is clear enough: after the catastrophe of the early 1930s, federal deposit insurance becomes part of the post-crisis architecture, supplementing central banking and helping to prevent the recurrence of classic retail bank runs on the same scale.

From this point forward, the mission of the central bank broadens. It is no longer only the institution that rescues banks in a panic. It becomes the institution that tries to stabilize the whole economy. When the economy overheats and inflation threatens, it raises rates and cools things down. When the economy weakens and unemployment rises, it cuts rates and encourages borrowing and spending.

This is where William McChesney Martin's phrase enters: the job of the central banker is to take away the punch bowl just as the party gets going. The lecture also recalls Charles Amos Dice's analogy. The Federal Reserve is like the regulator on a steam engine: not the source of power, but the device that prevents the system from running too fast.

The lecture then widens the lens geographically. The European Central Bank is mentioned as a comparatively recent example of central-banking invention. The Maastricht treaty laid the groundwork in 1992; the ECB was founded in 1998; the euro regime began in 1999, with physical currency coming later. The details are less important here than the lesson: central banking continues to evolve, and the same institutional logic can now govern a multinational currency area.

This is also where the lecture stresses independence. A central bank, on this view, should not be too easy for ordinary politics to push around, because governments under pressure may be tempted to inflate. Long terms for governors, insulation from dismissal, and the cultivation of a stable-currency ethos are not decorative constitutional features. They are part of the mechanism by which monetary credibility is built.

Only after that institutional background does the lecture narrow to current operating tools. The Federal Open Market Committee sets a target range for the federal funds rate, the overnight unsecured rate on loans between banks and certain other financial institutions. At the lecture date, the numbers are:

\begin{align}
0 \le i_{ff} \le 0.25\%, \qquad i_{ff} \approx 0.13\%.
\end{align}

The lecture's policy interpretation is straightforward. With unemployment high and the economy weak, the Fed is pushing rates essentially to zero. In the lecture's simplified presentation, zero is treated as the practical floor.

A second modern rate then appears. After the 2008 Emergency Economic Stabilization Act, the Federal Reserve begins paying interest on reserve balances:

\begin{equation}
i_{or} = 0.25\%.
\end{equation}

This is not presented as a minor technical adjustment. It is presented as a regime change. The old operating method worked mainly through open-market operations in Treasury bills and the indirect targeting of the fed funds rate. Interest on reserves gives the central bank a new tool, and therefore a new monetary-policy geometry.

\subsection{Question \& Answer}

\paragraph{Question.} Why can the interest on reserves sit above the federal funds rate without shutting the market down completely?

\paragraph{Answer.} Because the set of institutions active in the fed funds market is not identical to the set of institutions eligible to earn interest on reserves. If a member bank can earn
\[
i_{or} > i_{ff},
\]
it will usually prefer to leave money at the Federal Reserve rather than lend it overnight in fed funds. The lecture's answer is that some lenders in the fed funds market, notably government-sponsored enterprises such as Fannie Mae and Freddie Mac, are not eligible for interest on reserves. They therefore continue to lend there, and the two rates need not coincide.

\section{Reserve Requirements Under Regulation D}

At this point the lecture deliberately narrows again. We have been discussing the modern central bank as a macroeconomic regulator, but now the lecturer says, in effect, let us talk a little more about reserve requirements and let us distinguish them carefully from capital requirements. That return matters because reserve requirements once sat at the center of monetary theory.

The lecture also reminds us that reserve requirements and capital requirements did not begin with the Federal Reserve. State banking regulators in the United States were already using both kinds of requirements before 1913. The Federal Reserve later inherited and centralized those functions.

Reserve requirements are stated as a function of liabilities. In the lecture's language, reserves are currency in the vault plus deposits at the Federal Reserve. In cleaned notation we may write
\begin{equation}
R = \text{currency} + \text{deposits at the Federal Reserve}.
\end{equation}
This is the lecture's simplified notion of high-powered money from the viewpoint of the banking system.

Under Regulation D, the lecture-time example is:

\begin{figure}
\centering
\includegraphics[width=0.82\textwidth]{lecture_18_figure_02.png}
\caption{Reserve requirements under Regulation D. The board establishes the categories and the regulatory frame before the contrast is sharpened numerically.}
\end{figure}

\begin{figure}
\centering
\includegraphics[width=0.82\textwidth]{lecture_18_figure_03.png}
\caption{Transaction accounts versus time deposits. The later board view makes the \(10\%\) versus \(0\) comparison much clearer.}
\end{figure}

Near the second board, the clean mathematical reconstruction is:

\begin{align}
\text{Transaction accounts:} &\quad 10\%, \\
\text{Time deposits:} &\quad 0.
\end{align}

Equivalently,
\begin{align}
rr_{\text{transaction accounts}} &= 10\%, \\
rr_{\text{time deposits}} &= 0.
\end{align}

The chalkboard evidence supports exactly this side-by-side classification. It does not support more than that. The \(10\%\) figure and the zero on time deposits are visible; the fuller verbal explanation comes from the transcript. Transaction accounts are the immediately spendable, immediately withdrawable liabilities of the bank. Time deposits, in the lecture's framing, are liabilities on which payment can be delayed, so they do not present the same run problem.

\begin{remark}
The \(10\%\) and \(0\) figures here are lecture-time Regulation D values and should be read as the working example of the lecture, not as timeless constants.
\end{remark}

\subsection{Question \& Answer}

\paragraph{Question.} Why are transaction accounts subject to reserve requirements while time deposits are not?

\paragraph{Answer.} Because reserve requirements are aimed at the run problem. A transaction account holder can demand payment immediately. If many such depositors show up at once, the bank must have something ready to meet them. A time depositor does not have the same contractual right to immediate payment; the bank can stall, delay, or wait until the term structure of the liability allows repayment. So the reserve rule is attached to the liability that can run.

Once that distinction is established, the lecture reconstructs the older money-multiplier story. If the reserve ratio is \(rr\), and if banks hold only required reserves and no excess reserves, then deposits expand until reserves are exactly the required fraction of deposits. In its simplest lecture form,
\begin{equation}
m \approx \frac{1}{rr}.
\end{equation}
With \(rr = 0.10\), this becomes
\begin{equation}
m \approx 10.
\end{equation}

The idea is that if banks are willing to sit exactly on the legal minimum, then reserves pin down the deposit base. Under that old theory, reserve requirements are close to being the central lever of money-supply determination.

\begin{remark}
The lecture is explicit that this formula is an oversimplification. It is reconstructed here because it captures the older logic of reserve-based banking, not because it is an exact identity.
\end{remark}

The lecture then immediately limits the old theory. The world has changed. Before the recent crisis, banks did not want to hold excess reserves because reserves earned no interest, so reserve balances sat not far above the minimum. But once the Federal Reserve began paying interest on reserves, banks became willing, even eager, to hold very large excess balances. At the lecture date, the quantity is described as enormous:
\begin{equation}
\text{Excess reserves} \gtrsim 1.2 \text{ trillion dollars}.
\end{equation}
At that point reserve requirements are no longer the binding center of the system for most banks. The lecture therefore uses the reserve-requirement section to recover an older theory and then to say, bluntly, that this is not the world we are living in now.

\section{Capital Requirements, Risk Weights, and the Basel III Shift}

This brings us, in the lecture's own rhythm, to capital requirements. Reserve requirements are defined off the liability side of the balance sheet. Capital requirements are defined off the asset side, after assets have been adjusted for risk.

In cleaned notation, the lecture's risk-weighted-asset idea may be written as
\begin{equation}
\text{RWA} = \sum_i w_i A_i,
\end{equation}
where \(A_i\) denotes an asset category and \(w_i\) its regulatory risk weight. The later toy example uses two especially simple weights:
\begin{equation}
w_{\text{cash}} = 0, \qquad w_{\text{corporate loans}} = 1.
\end{equation}

The Basel III common-equity arithmetic is then summarized in the lecture as follows:
\begin{align}
\text{CET1 requirement} &= 4.5\% \times \text{RWA}, \\
\text{capital conservation buffer} &= 2.5\% \times \text{RWA}.
\end{align}
So the ordinary simplified common-equity target becomes
\begin{equation}
7\% \times \text{RWA}.
\end{equation}
And the lecture adds that regulators may impose a countercyclical buffer of up to
\begin{equation}
2.5\% \times \text{RWA},
\end{equation}
raising the total common-equity requirement to
\begin{equation}
9.5\% \times \text{RWA}.
\end{equation}

\begin{remark}
The spoken lecture is loose at moments in the verbal arithmetic of Basel III. The cleaned reconstruction above preserves the intended structure: \(4.5\%\) minimum common equity, \(2.5\%\) conservation buffer, and up to another \(2.5\%\) countercyclical addition.
\end{remark}

\subsection{Question \& Answer}

\paragraph{Question.} What is the difference between reserve requirements and capital requirements?

\paragraph{Answer.} Reserve requirements are a liquidity rule. They ask whether the bank has enough immediately available reserves against liabilities that may run. Capital requirements are a solvency rule. They ask whether the bank has enough loss-absorbing common equity relative to the riskiness of its assets. A bank can satisfy the first and fail the second; the toy example later in the lecture is built precisely to show that possibility.

The lecture then adds a modern complication: how do we measure risk? For many years regulators leaned heavily on the ratings issued by recognized credit-rating agencies. This is where the acronym NRSRO appears:
\begin{equation}
\text{NRSRO} = \text{Nationally Recognized Statistical Rating Organization}.
\end{equation}
Moody's and Standard \& Poor's are the main examples in the lecture. A large regulatory structure grew up around the assumption that such agencies could map the risk of securities reliably enough for capital rules to lean on them.

The crisis damaged that confidence badly. Securities that carried top ratings performed disastrously. So the lecture presents Dodd--Frank as a sharp reaction against formal regulatory dependence on such ratings. The rhetoric in the transcript is stronger than the careful legal formulation; the right way to state the point here is that Dodd--Frank pushes regulators away from using ratings as automatic regulatory truth. Basel III, however, still refers to ratings in many places. That creates a tension for American implementation. The lecture's deeper point is that capital regulation is only as good as the risk measurement built into it, and risk measurement turned out to be far less secure than people had imagined.

\section{A Toy Bank Balance Sheet and the Crisis Logic of Regulation}

The lecture's main mathematical construction is not a theorem. It is a toy bank. We found a bank, open the doors, accept a deposit, issue equity, make loans, absorb losses, and then ask which regulations still bind. This step-by-step construction is the real analytical spine of the chapter.

We begin with the balance-sheet identity
\begin{equation}
A = L + E_{\text{common}},
\end{equation}
where \(A\) is total assets, \(L\) non-equity liabilities, and \(E_{\text{common}}\) common equity.

Suppose someone walks in and deposits \$100 in cash. The first balance sheet is
\begin{align}
A &= 100 \text{ cash}, \\
L &= 100 \text{ transaction account}, \\
E_{\text{common}} &= 0.
\end{align}

If the deposit is a transaction account, required reserves are
\begin{equation}
0.10 \times 100 = 10.
\end{equation}
So from the viewpoint of reserve requirements the bank is fine. It has \$100 in the vault against a \$10 requirement.

Now look at capital. Because cash has zero risk weight,
\begin{equation}
\text{RWA} = 0.
\end{equation}
At this knife-edge starting point, the ratio rule itself is vacuous:
\begin{equation}
4.5\% \times \text{RWA} = 0.
\end{equation}
But the lecture uses this opening to make a practical point. A bank with no common equity at all is not the kind of institution regulators want to rely on. So the story immediately injects capital by issuing shares.

Suppose the bank issues \$20 of common equity. Then the balance sheet becomes
\begin{align}
A &= 120 \text{ cash}, \\
L &= 100 \text{ transaction account}, \\
E_{\text{common}} &= 20.
\end{align}

The lecture pauses here to emphasize why common equity is different from deposits. The transaction-account holder can come to the window and demand money immediately. The shareholder cannot. Shareholders own the residual claim on the bank; they do not have the right to run it at the window. That is why common equity is the loss-absorbing liability.

At this stage the bank is safe but unprofitable. All assets are still cash. So the bank now does what banks do: it lends. It converts \$100 of cash into corporate loans and keeps \$20 in cash:
\begin{align}
A &= 100 \text{ corporate loans} + 20 \text{ cash}, \\
L &= 100 \text{ transaction account}, \\
E_{\text{common}} &= 20.
\end{align}

Using the lecture's simplified risk weights,
\begin{equation}
\text{RWA} = 100.
\end{equation}
Under the simplified Basel III target,
\begin{equation}
0.07 \times \text{RWA} = 7.
\end{equation}
So the bank is comfortably above its capital requirement: it has \$20 in common equity against a \$7 requirement.

It is also still above its reserve requirement. The transaction account is still \$100, so the reserve requirement is still \$10, and the bank is holding \$20 in cash. At this moment the toy bank satisfies both regulatory systems at once.

Now comes the crisis. Suppose \(20\%\) of the corporate-loan portfolio defaults, so the bank writes the loans down from \$100 to \$80. Then the assets become
\begin{align}
A &= 80 \text{ corporate loans} + 20 \text{ cash} = 100.
\end{align}
The deposit liability is still \$100. Therefore the only thing that can give way is common equity:
\begin{align}
L &= 100 \text{ transaction account}, \\
E_{\text{common}} &\to 0.
\end{align}

This is the crucial contrast in the lecture. Reserve requirements are still satisfied:
\begin{equation}
\text{required reserves} = 10, \qquad \text{cash held} = 20.
\end{equation}
But capital requirements are violated. With the same simplified risk weights, the bank now has
\begin{equation}
\text{RWA} = 80,
\end{equation}
so even the \(7\%\) common-equity target requires
\begin{equation}
0.07 \times 80 = 5.6,
\end{equation}
and the bank has zero. Reserve adequacy and capital adequacy have come apart.

This is exactly what the lecture wants us to see. Reserve requirements protect against immediate withdrawal pressure. Capital requirements protect against asset losses. A bank can still have enough cash to meet withdrawals and yet have no loss-absorbing equity left.

The lecture then stages a false repair and corrects it. At first one might think the bank can solve the problem by selling loans and holding more cash. Suppose it sells \$20 of corporate loans and receives cash at book value. Then the balance sheet would be
\begin{align}
A &= 60 \text{ corporate loans} + 40 \text{ cash} = 100, \\
L &= 100 \text{ transaction account}, \\
E_{\text{common}} &= 0.
\end{align}
Risk-weighted assets would fall to \(60\), but common equity would still be zero. So this does not solve the problem:
\begin{equation}
0 < 0.07 \times 60 = 4.2.
\end{equation}
The lecturer explicitly corrects himself on this point. Once common equity has been driven all the way to zero, merely reshuffling the asset side does not recreate equity. If equity had fallen only part of the way, then selling risky assets could help by lowering \(\text{RWA}\). But once equity is exhausted, the only genuine repair is to issue new shares.

That repair is easy to write on the board and hard to carry out in a crisis. The previous shareholders have just been wiped out. New investors know this. The lecture therefore widens the story from one bank to all banks. If a whole system is hit at once, every bank tries to do the same two things: raise new equity or sell assets. Both actions are hardest precisely when the crisis is most severe. That is how a prudential rule meant to stabilize the system can become procyclical.

\subsection{Question \& Answer}

\paragraph{Question.} Why can capital requirements become procyclical and dangerous in a crisis even if they are meant to make banks safer?

\paragraph{Answer.} Because the requirement bites hardest after losses. Once equity has been eroded, a bank must either raise new equity or shrink risky assets. But in a crisis new equity is hard to sell, and asset sales become fire sales because every damaged bank is trying to sell at the same time. Prices fall, balance sheets weaken further, and the system contracts together. The rule is individually sensible, but systemically it can intensify the downturn unless large buffers are built before the crisis.

That is the final turn of the lecture. Basel III's countercyclical buffer is an attempt to move the adjustment earlier in time, forcing thicker equity cushions when a bubble is building rather than waiting until the crash. At the same time, Dodd--Frank constrains the Federal Reserve's discretionary rescue power by pushing it toward a more uniform discount-window style of support rather than ad hoc bailouts of selected institutions. The system was saved in the recent crisis by lender-of-last-resort action, but the lecture ends by insisting that the design problem is not solved. The institutions are still evolving.

\section{Summary}

The lecture begins by redefining the subject. A central bank is not just a government office with technical powers over money. It is a copied financial invention. The Bank of England supplies the original durable model: reserves concentrated at a dominant institution, support in trouble, and the emergence of a lender of last resort. The United States then spends a century moving unevenly toward the same logic through Suffolk, national banking, and finally the Federal Reserve.

The second half of the lecture reconstructs the old reserve-based theory only in order to historicize it. Under Regulation D, transaction accounts carried a \(10\%\) reserve requirement and time deposits a zero requirement; this supported the familiar oversimplified multiplier \(m \approx 1/rr\). But once banks came to hold massive excess reserves, reserve requirements ceased to be the main live constraint. Capital requirements moved to the center instead.

The toy bank makes that change precise. A bank can satisfy reserve requirements and still fail capital requirements after losses. That is why modern prudential regulation focuses on common equity relative to risk-weighted assets. But the lecture's final point is uneasy rather than triumphant: a system that forces banks to repair capital in the middle of a crisis may itself deepen the crisis. Central banks saved the system last time. Whether the new regulatory design can prevent the same pattern from recurring remains an open question.